\cotefile{93}
\scanpage{61}
\mspage{8.}
% archive sheet number: 60

\diagram{\(\begin{array}{@{:~}l@{~:~}c@{~:~}c@{~:~}c@{~:~}c@{~:~}c@{~:~}c}
\hline
 & 1971 & 1972 & 1973 & 1974 & 1975 & 1976\\ \hline
\begin{array}{l}\text{Fonctionnement}\\ \text{(MF) Chap. 36.15}\end{array}
 & 190 & 206 & 223 & 249 & 257 & 275\\ \hline
\begin{array}{l}\text{Equipement (MF)}\\ \text{Chap. 56-10 (50)}\end{array}
 & 87 & 92 & 93 & 93 & 104 & 79\\ \hline
\text{TOTAL} & 277 & 298 & 316 & 342 & 361 & 3\deleted{\unsure{9|5}}\handwritten{5}4\\ \hline
\end{array}\)}

La quasi totalité du crédit de fonctionnement est réparti en tenant compte de critères quantitatifs :

- nombre d'enseignant chercheurs de l'établissement

- nombre de thèses soutenues

- nombre d'association réalisées avec d'autres établissements de recherche

- investissements réalisés au cours des six dernières années au profit des formations de recherche.

- nouvelles actions originales de l'établissement.

Cette répartition est effectuée avant le début de l'année civile. La totalité du budget 1976 est donc actuellement répartie.

Le budget d'équipement est réparti par la D.E.S.U.R. sous la forme de deux grandes masses :

- des constructions spécifiques pour la recherche, et des investissements de matériels de grande importance (ordinateurs, matériels communs à plusieurs établissements) ; ces opérations sont pilotés par la DESUR en accord avec les établissements.

- du matériel "moyen" pour la recherche, ces appareils sont attribués à l'établissement sur sa proposition, après examen et approbation par les experts de la DESUR.
